\begin{abstract}
The $XYZ^2$ hexagonal stabilizer code concatenates a YZZY surface code with two-qubit phase-flip codes on the links of a honeycomb lattice. We study the preparation of its logical state $\ket{\bar{+}}$ by measurement, decoding and recovery (MDR) under circuit-level noise. We construct a fault-tolerant MDR protocol in which the data qubits start in a product state chosen in the CSS frame of the code, all generators are measured with an interleaved depth-six schedule selected by an exhaustive search, and the full space-time syndrome is decoded before a Pauli frame update. The protocol has circuit-level fault distance $d+1$, which we verify exactly for $d=3$ and $5$. With distances up to 19, the threshold of matching under circuit-level depolarizing noise falls from $1.8\%$ for a single round of measurements to $0.43\%$ when the number of rounds grows with the distance. We introduce two decoders. Hierarchical correlated matching uses the gauge checks as a lower decoding level and raises the threshold of matching by $22\%$ for depolarizing noise and by $53\%$ for $Z$-biased noise with bias 100, at about four times the cost of matching. The coset free-energy decoder compares the two logical classes directly. It combines a two-class ordered-statistics search with local moves between equivalent errors and an entropy correction. It fails on as few shots as the search-based Tesseract decoder at $d=5$, and under strongly biased noise on $28\%$ fewer shots than belief-matching at $d=7$. For hardware with native pair measurements we give an extraction with cat-state ancillas, which has fault distance $d$ and a threshold of $0.28\%$ under the EM3 noise model. We describe the Quantinuum Helios and H2 processors by one-parameter noise models with the error ratios of their data sheets and check them against independent measurements and the vendor's emulator. At the Helios operating point a single round prepares $\ket{\bar{+}}$ at $d=5$ with a logical error rate of $8\times10^{-5}$, six times below the physical SPAM error. Measurement crosstalk, which grows with the number of measured ancillas, moves the crossing points of consecutive distances to smaller error rates as the code grows.
\end{abstract}
